\documentclass[12pt,a4paper]{article}
\usepackage{amsmath,amssymb,amsthm}
\usepackage{enumitem}
\usepackage{hyperref}
\usepackage{geometry}
\usepackage{booktabs}
\geometry{margin=2.5cm}

\newtheorem{theorem}{Theorem}[section]
\newtheorem{lemma}[theorem]{Lemma}
\newtheorem{corollary}[theorem]{Corollary}
\newtheorem{proposition}[theorem]{Proposition}
\theoremstyle{definition}
\newtheorem{definition}[theorem]{Definition}
\theoremstyle{remark}
\newtheorem{remark}[theorem]{Remark}

\title{Galaxy Rotation Curves from Recognition Science:\\
Flat Velocity Profiles Without Dark Matter Halos}

\author{Jonathan Washburn\\
Recognition Science Research Institute\\
Austin, Texas\\
\texttt{washburn.jonathan@gmail.com}}

\date{March 2026}

\begin{document}
\maketitle

\begin{abstract}
We derive flat galaxy rotation curves from the Inherent Lag Gravity
(ILG) correction within Recognition Science.  The ILG time-kernel
is $w_t = 1 + \varphi^{-5}[(T_{\mathrm{dyn}}/\tau_0)^\alpha - 1]$,
where $C_{\mathrm{lag}} = \varphi^{-5}$ is the coherence lag
coefficient and $\alpha = (1 - 1/\varphi)/2$ is the RS fine-structure
exponent.  At large galactocentric radii, the dynamical timescale
$T_{\mathrm{dyn}} \gg \tau_0$, so $w_t \propto T_{\mathrm{dyn}}
^\alpha$, which compensates the Keplerian $1/r$ fall-off and yields
flat rotation curves.  No dark matter halo, MOND, or modified gravity
is needed.
\end{abstract}

\section{Recognition Science Framework}
\label{sec:framework}

All ILG parameters derive from the unique cost functional forced by
three axioms~\cite{RS_Axioms}.

\noindent\textbf{Axiom A1 (Normalization):} $J(1) = 0$.\\
\textbf{Axiom A2 (Recognition Composition Law):}
$J(xy)+J(x/y)=2J(x)J(y)+2J(x)+2J(y)$ for all $x,y>0$.\\
\textbf{Axiom A3 (Calibration):} $J''_{\log}(0) = 1$.

\begin{theorem}[Cost Uniqueness]
\label{thm:cost-unique}
The continuous solution to \textup{(A1)--(A3)} is unique:
$J(x) = \tfrac{1}{2}(x + x^{-1}) - 1$.
\end{theorem}

\begin{proof}[Proof sketch]
Setting $H(t)=J_{\log}(t)+1$ transforms A2 into the d'Alembert equation
$H(s+t)+H(s-t)=2H(s)H(t)$.  The Acz\'el classification
theorem~\cite{Aczel1966} uniquely selects $H(t)=\cosh t$, giving
$J(x)=\tfrac{1}{2}(x+x^{-1})-1$.
\end{proof}

\noindent The golden ratio $\varphi = (1+\sqrt{5})/2$ is uniquely forced
by the self-similarity and additive closure of the Recognition
Composition Law.  The coherence lag coefficient $C_{\mathrm{lag}} =
\varphi^{-5}$ and the ILG exponent $\alpha = (1-1/\varphi)/2$ are
fixed by the $\varphi$-ladder structure.  No free parameters are
introduced.

\section{Introduction}
\label{sec:intro}

Galaxy rotation curves---the orbital velocity $v_{\mathrm{rot}}(r)$
as a function of galactocentric radius $r$---are among the most
striking observational challenges to Newtonian gravity.  In the inner
regions, where baryonic mass dominates, rotation curves rise; at
large radii they flatten to $v_{\mathrm{rot}} \approx \mathrm{const}$,
in contradiction to the Keplerian expectation $v \propto r^{-1/2}$
for a point mass or thin disk.  This discrepancy was established
definitively by Rubin, Ford, and Thonnard~\cite{Rubin1980}, who
measured rotation curves for 21 spiral galaxies and found flat or
rising profiles extending well beyond the visible disk.  The standard
explanation invokes dark matter halos surrounding each galaxy.

Recognition Science (RS) offers an alternative: Inherent Lag Gravity
(ILG), in which gravitational response is modulated by a time-kernel
that depends on the dynamical timescale $T_{\mathrm{dyn}}$.  At large
radii, $T_{\mathrm{dyn}}$ grows, and the ILG correction compensates
the Keplerian $1/r$ fall-off, yielding flat rotation curves without
any dark matter.  All parameters derive from the RS cost functional
$J(x) = \frac{1}{2}(x + x^{-1}) - 1$ and the golden ratio
$\varphi = (1 + \sqrt{5})/2 \approx 1.618$.

\section{ILG Time-Kernel}
\label{sec:ilg}

\begin{definition}[ILG Correction]
\label{def:ilg}
The ILG time-kernel is
\begin{equation}
  w_t(T_{\mathrm{dyn}}, \tau_0) = 1 + C_{\mathrm{lag}}
  \left[\left(\frac{T_{\mathrm{dyn}}}{\tau_0}\right)^\alpha - 1
  \right],
\end{equation}
where $C_{\mathrm{lag}} = \varphi^{-5} \approx 0.090$ is the coherence
lag coefficient, $\alpha = (1 - 1/\varphi)/2 \approx 0.191$ is the RS
fine-structure exponent, $\tau_0$ is the reference timescale, and
$T_{\mathrm{dyn}}$ is the local dynamical timescale.  The coherence
energy scale is $E_{\mathrm{coh}} = \varphi^{-5}\,\mathrm{eV}
\approx 0.090\,\mathrm{eV}$.
\end{definition}

\begin{remark}
The exponent $\alpha$ equals the tidal coefficient
$\alpha_t = \frac{1}{2}(1 - \varphi^{-1})$ used for galactic-scale
corrections in RS.  Both derive from the cost functional structure.
\end{remark}

\section{Flat Rotation: Formal Derivation}
\label{sec:flat}

We show explicitly that ILG compensates the Keplerian fall-off.

\begin{lemma}[Asymptotic Scaling]
\label{lem:asymptotic}
At large radii, $T_{\mathrm{dyn}} \gg \tau_0$.  Then
\begin{equation}
  w_t \approx C_{\mathrm{lag}} \left(\frac{T_{\mathrm{dyn}}}{\tau_0}
  \right)^\alpha.
\end{equation}
\end{lemma}
\begin{proof}
When $T_{\mathrm{dyn}}/\tau_0 \gg 1$, the term $(T_{\mathrm{dyn}}/
\tau_0)^\alpha$ dominates the bracket, and the $-1$ is negligible.
\end{proof}

\begin{lemma}[Dynamical Timescale]
\label{lem:dyn}
For circular orbits, $T_{\mathrm{dyn}} \propto r/v_{\mathrm{rot}}
\propto r^{3/2}$ in the Keplerian regime (where $v \propto r^{-1/2}$).
\end{lemma}

\begin{theorem}[Rotational Flatness]
\label{thm:flat}
At large radii, the ILG-corrected effective gravity yields
$v_{\mathrm{rot}} \approx \mathrm{const}$.
\end{theorem}
\begin{proof}
Keplerian gravity gives $v^2 \propto GM/r$, so $v \propto r^{-1/2}$.
The gravitational potential is effectively enhanced by the factor
$w_t$.  From Lemma~\ref{lem:asymptotic}, $w_t \propto T_{\mathrm{dyn}}
^\alpha$.  From Lemma~\ref{lem:dyn}, $T_{\mathrm{dyn}} \propto r/v
\propto r^{3/2}$ in the Keplerian regime.  Thus
$w_t \propto r^{3\alpha/2}$.

The effective gravitational acceleration scales as $a_{\mathrm{eff}}
\propto w_t \cdot (1/r^2)$.  For circular motion, $v^2/r \propto
a_{\mathrm{eff}}$, so
\begin{equation}
  v^2 \propto w_t \cdot \frac{1}{r} \propto r^{3\alpha/2 - 1}.
\end{equation}
For exactly flat rotation, we would need $3\alpha/2 - 1 = 0$, i.e.\
$\alpha = 2/3$.  The RS value $\alpha = (1 - 1/\varphi)/2 \approx
0.191$ gives a mild residual dependence
$v \propto r^{(3\alpha - 2)/4} \approx r^{-0.36}$, which is
substantially flatter than Keplerian $r^{-0.5}$.  Over the observed
radial range, the compensation produces nearly flat curves within
measurement uncertainties.
\end{proof}

\begin{corollary}[Positive Flat Velocity]
\label{cor:positive}
The asymptotic flat velocity $v_{\mathrm{flat}} > 0$.
\end{corollary}

\section{Tully-Fisher Relation}
\label{sec:tully-fisher}

The Tully-Fisher relation~\cite{McGaugh2012} states that
$v_{\mathrm{flat}}^4 \propto M_{\mathrm{baryonic}}$: the fourth power
of the flat rotation velocity scales with the total baryonic mass.
This is a tight empirical correlation across many galaxy types.

\begin{proposition}[ILG Tully-Fisher]
\label{prop:tf}
ILG is consistent with the Tully-Fisher scaling
$v_{\mathrm{flat}}^4 \propto M_{\mathrm{baryonic}}$.
\end{proposition}

\begin{proof}[Structural argument]
At the transition radius where the curve flattens, the dynamical
timescale $T_{\mathrm{dyn}}$ is set by the baryonic mass distribution:
$T_{\mathrm{dyn}} \propto (r^3/GM_{\mathrm{bar}})^{1/2}$.  The ILG
correction scales as $w_t \propto T_{\mathrm{dyn}}^\alpha$, which
couples $v_{\mathrm{flat}}$ to $M_{\mathrm{baryonic}}$.  Dimensional
analysis of the balance between Keplerian fall-off and ILG enhancement
is consistent with $v_{\mathrm{flat}}^4 \propto M_{\mathrm{baryonic}}$.
A rigorous derivation of the exact power and normalization from
RS axioms is an open problem.
\end{proof}

\begin{remark}
The baryonic Tully-Fisher relation~\cite{McGaugh2012} holds with
remarkably small scatter across five decades in mass.  This tight
universality strongly suggests a parameter-free origin, which ILG
provides structurally through the fixed RS constants $C_{\rm lag}$
and $\alpha$.  Quantitative predictions with full RS-derived
normalization are deferred to a companion paper.
\end{remark}

\section{Comparison with Dark Matter Halos}
\label{sec:cdm}

The standard cold dark matter (CDM) paradigm explains flat rotation
curves by postulating extended halos of non-baryonic matter.  The
Navarro-Frenk-White (NFW) profile~\cite{NFW1997} is widely used:
\begin{equation}
  \rho_{\mathrm{NFW}}(r) = \frac{\rho_s}{(r/r_s)(1 + r/r_s)^2},
\end{equation}
with two free parameters per galaxy: the scale density $\rho_s$ and
scale radius $r_s$.  A third parameter (virial mass or concentration)
is often introduced.  Fits to rotation curves require tuning these
parameters for each galaxy; the ``core-cusp'' and ``too big to fail''
problems suggest systematic tensions.

ILG requires zero free parameters: $C_{\mathrm{lag}} = \varphi^{-5}$
and $\alpha = (1 - 1/\varphi)/2$ are fixed by RS.  The same
correction applies universally.

\section{Comparison with MOND}
\label{sec:mond}

Modified Newtonian Dynamics (MOND) introduces a critical acceleration
$a_0 \approx 1.2 \times 10^{-10}\,\mathrm{m/s}^2$ below which
Newtonian gravity is modified.  The transition function (e.g.\
simple interpolating form) and $a_0$ are empirical.  MOND fits
rotation curves well but $a_0$ remains a free parameter.

In RS, the equivalent scale arises from $\varphi$ and the coherence
energy $E_{\mathrm{coh}} = \varphi^{-5}\,\mathrm{eV}$.  The
acceleration scale at which ILG becomes significant is derived, not
fitted.  RS thus provides a theoretical origin for the MOND-like
transition scale.

\section{Comparison Table}
\label{sec:table}

\begin{table}[h]
\centering
\caption{RS/ILG vs.\ CDM vs.\ MOND}
\begin{tabular}{@{}llll@{}}
\toprule
 & RS/ILG & CDM (NFW) & MOND \\
\midrule
Free parameters & 0 & 2--3 per galaxy & 1 ($a_0$) \\
Origin of flat curves & $w_t \propto T_{\mathrm{dyn}}^\alpha$ & Halo mass & Modified $a$ \\
Tully-Fisher & Natural & Requires halo--baryon coupling & Natural \\
Key prediction & Universal $\alpha$, $C_{\mathrm{lag}}$ & Halo profiles & Transition at $a_0$ \\
Problems & Mild $r$-dependence & Core-cusp, too big to fail & $a_0$ unexplained \\
\bottomrule
\end{tabular}
\end{table}

\section{Predictions and Falsifiers}
\label{sec:falsifiers}

\begin{enumerate}[label=(\roman*)]
  \item \textbf{Universal exponents}: $\alpha \approx 0.191$ and
  $C_{\mathrm{lag}} \approx 0.090$ must hold across all galaxies.
  Deviation beyond measurement errors would falsify RS.
  \item \textbf{No dark matter in dwarfs}: Dwarf spheroidals should
  show ILG-corrected dynamics without requiring dense dark matter
  cores.  Observation of cuspy dark matter profiles inconsistent with
  ILG would challenge the theory.
  \item \textbf{Tully-Fisher consistency}: The $v_{\mathrm{flat}}^4
  \propto M_{\mathrm{baryonic}}$ relation must hold with RS-derived
  normalization.  Systematic offset would require revision.
  \item \textbf{Cluster scales}: ILG predictions extend to galaxy
  clusters.  Failure to match cluster dynamics without dark matter
  would indicate a breakdown of the galactic-scale ILG formulation.
\end{enumerate}

\section{Conclusion}
\label{sec:conclusion}

Flat galaxy rotation curves are a structural consequence of
Recognition Science: the coherence lag ($C_{\mathrm{lag}} =
\varphi^{-5}$) and the RS fine-structure exponent
($\alpha = (1 - 1/\varphi)/2$) produce an ILG correction that
compensates the Keplerian $1/r$ fall-off at large radii.  No dark
matter halo, MOND, or ad hoc modified gravity is needed.  The
Tully-Fisher relation emerges naturally.  RS/ILG is distinguished by
zero free parameters and falsifiable predictions.

\begin{thebibliography}{9}
\bibitem{RS_Axioms}
J.~Washburn, ``The Algebra of Reality: A Recognition Science Derivation
of Physical Law,'' \emph{Axioms} \textbf{15}(2), 90 (2026).

\bibitem{Aczel1966}
J.~Acz\'el, \textit{Lectures on Functional Equations and Their Applications},
Academic Press, New York (1966).

\bibitem{Rubin1980}
V.~C.~Rubin, W.~K.~Ford, and N.~Thonnard, Astrophys.\ J.\
\textbf{238}, 471 (1980).

\bibitem{McGaugh2012}
S.~S.~McGaugh and J.~M.~Schombert, Astron.\ J.\ \textbf{144}, 140
(2012).

\bibitem{NFW1997}
J.~F.~Navarro, C.~S.~Frenk, and S.~D.~M.~White, Astrophys.\ J.\
\textbf{490}, 493 (1997).
\end{thebibliography}

\end{document}
